% Management in Global Markets... Business Plan (Tierra Bar)
% Different format from CaféOlé. Orange theme.
% Compile: pdflatex business-plan-tierra-bar.tex
% Font: Comic Neue Angular (handwritten feel, different from CafeOle)

\documentclass[11pt,a4paper]{article}
\usepackage[utf8]{inputenc}
\usepackage[T1]{fontenc}
\usepackage[default,angular]{comicneue}
\usepackage[margin=1.8cm, top=2.2cm, bottom=2cm]{geometry}
\usepackage{ragged2e}
\RaggedRight
\usepackage{parskip}
\setlength{\parskip}{0.6em}
\usepackage{enumitem}
\usepackage[table]{xcolor}
\usepackage{booktabs}
\usepackage{titlesec}
\usepackage{fancyhdr}
\usepackage{hyperref}
\usepackage{microtype}
\hyphenation{study-abroad}
\definecolor{tierraorange}{RGB}{230,115,40}
\definecolor{darkorange}{RGB}{180,85,25}
\definecolor{lightorange}{RGB}{255,240,220}
\definecolor{charcoal}{RGB}{51,51,51}

% Section style: orange header (notes feel)
\titleformat{\section}[block]{\color{tierraorange}\normalfont\large\bfseries}{\thesection.}{0.5em}{}[\vspace{-0.2em}\textcolor{tierraorange}{\rule{\textwidth}{1.2pt}}\vspace{0.4em}]
\titlespacing*{\section}{0pt}{1.2em}{0.6em}

\pagestyle{fancy}
\fancyhf{}
\fancyhead[C]{\small\textcolor{tierraorange}{TIERRA BAR}}
\fancyfoot[R]{\color{tierraorange}\thepage}
\renewcommand{\headrulewidth}{0.4pt}
\renewcommand{\footrulewidth}{0pt}

\hypersetup{
  colorlinks=true,
  linkcolor=tierraorange,
  urlcolor=tierraorange,
  pdftitle={Tierra Bar Business Plan},
  pdfauthor={Ronan Smith, Gil Arvin, Daniel Tryder, Jacob Peterson},
  pdfsubject={Global Management}
}

\begin{document}

% ===== FRONT PAGE (different layout: left-aligned, minimal) =====
\thispagestyle{empty}
\vspace*{2cm}
\noindent
\fcolorbox{tierraorange}{lightorange}{\parbox{0.95\textwidth}{
\vspace{0.8em}
{\large\bfseries\color{darkorange} TIERRA BAR}\\[0.2em]
{\normalsize Business Plan}\\[0.2em]
{\normalsize\itshape Management in Global Markets... applying class frameworks.}\\[0.8em]
{\normalsize Ronan Smith \textbullet\ Gil Arvin \textbullet\ Daniel Tryder \textbullet\ Jacob Peterson}\\[0.3em]
{\small Prof. Geoffrey Ditta \textbullet\ March 9, 2026}\\[0.3em]
{\small \copyright\ 2026 Tierra Burrito Bar}\\[0.8em]
}}
\vfill
\noindent{\small Facultad de Geograf\'ia e Historia \textbullet\ Universidad Complutense de Madrid}
\newpage
\pagestyle{fancy}
\setcounter{page}{1}

% ===== 1. OVERVIEW =====
\section{Overview}\label{sec:overview}

\textbf{So here's the deal.} Tierra Bar (Tierra Burrito Bar)... Plaza de Isabel II, 5, Madrid Centro, near \'Opera. Build-your-own burritos, tacos, nachos. Margaritas, 16 craft taps. Closes 1:30 AM weekends. Glovo, Just Eat. 9/10, 3,500+ reviews. Madrid, Barcelona, Salamanca.

From class: a business plan explains the opportunity, identifies the market, and details how we pursue it (Harvard Business School definition).

\begin{itemize}[leftmargin=*,nosep]
\item \textbf{Goal (scope 4.1):} Position Tierra Bar as the late-night Mexican spot for students (Erasmus, study abroad)... affordable food, drinks, community.
\item \textbf{Plan:} 15\% student discount, Tierra Student Pass (20 euros/month: one free meal/week + drink discounts), reserved late-night tables, events (Taco Tuesday language exchange, Margarita Meetups).
\item \textbf{Competitive advantage (from class):} Few bars in Madrid combine late hours, craft beer, quality food, student offers, and community. We have the location and reviews. When students find community, they return. That increases profit and reputation.
\end{itemize}

\vspace{1em}

% ===== 2. EXECUTIVE SUMMARY =====
\section{Executive Summary}

\textbf{Alright, let me break this down.} Tierra Bar is a Mexican bar-restaurant in Madrid Centro. Current: tourists, professionals, locals. We're adding students. Erasmus and study abroad need affordable late-night options. Student Pass (20 euros/month) + 15\% discount + events = community. When students find community, they come back. That's our competitive advantage from class.

\textbf{Market gap:} Madrid has bars and Mexican spots, but almost none target students with discounts and events. We're filling it. ``Go where competitors are''... they validate the market. We differentiate.

\textbf{PESTLE (macro-environment from class):} Political... SL, Madrid licenses, alcohol license. Economic... tourism 11M, unemployment 9\% (lowest since 2008), student budgets. Social... demand for casual Mexican, late-night venues, veggie options. Tech... reviews, social media, delivery apps. Legal... AESAN, HACCP, minimum wage (1,134 euros full-time, 620 part-time). Environmental... sustainable packaging.

\textbf{SWOT (from class):} Strengths... location, food, craft beer, late hours, reviews. Weaknesses... no student programme yet. Opportunities... Complutense, UAM, Erasmus, tourism. Threats... Taco Bell, taquer\'ias, cost increases.

\textbf{Key takeaway:} Viable. Execute.

\vspace{1em}

% ===== 3. COMPANY PROFILE =====
\section{Company Profile}

\begin{tabular}{@{}ll@{}}
\textbf{Legal form} & Tierra Burrito Bar SL (Sociedad Limitada) \\
\textbf{Address} & Plaza de Isabel II, 5, 28013 Madrid \\
\textbf{Offer} & Mexican fast-casual (burritos, tacos, nachos, bowls), bar (margaritas, 16 craft taps), late hours, delivery \\
\textbf{Student programme} & 15\% discount, Student Pass 20 euros/month, events \\
\textbf{Mission} & Affordable Mexican food and a welcoming late-night space \\
\textbf{Values} & Fresh ingredients, generous portions, community \\
\textbf{Tagline} & Good food, good drinks, good people \\
\textbf{Course link} & Resources (kitchen, bar, space) + capabilities (events, partnerships) \\
\end{tabular}

\vspace{1em}

% ===== 4. BUSINESS MODEL CANVAS =====
\section{Business Model Canvas}\label{sec:bmc}

\textbf{From class:} The BMC has nine blocks. Here's ours:

\begin{center}
\small
\begin{tabular}{@{}p{5.2cm}p{5.2cm}@{}}
\toprule
\rowcolor{lightorange}
\textbf{Customer segments} & \textbf{Value proposition} \\
\midrule
Students (Erasmus, study abroad), young professionals, tourists, locals & Customisable Mexican food, craft beer, margaritas, late hours, student discounts, events \\
\addlinespace
\rowcolor{lightorange}
\textbf{Channels} & \textbf{Revenue streams} \\
\midrule
In-store, Glovo, Just Eat, tierraburritos.com, Instagram, TikTok & Food, bar, delivery, Student Pass (20 euros/month) \\
\addlinespace
\rowcolor{lightorange}
\textbf{Key resources} & \textbf{Key activities} \\
\midrule
Kitchen, bar, 16 taps, space, ingredients, staff & Food prep, bar service, events, university partnerships \\
\addlinespace
\rowcolor{lightorange}
\textbf{Key partnerships} & \textbf{Cost structure} \\
\midrule
Complutense, UAM, Glovo, Just Eat, beer suppliers & Rent 20--30\%, supplies 30--40\%, labour, delivery commission, marketing 2--5\% \\
\bottomrule
\end{tabular}
\end{center}

\vspace{1em}

% ===== 5. PESTLE =====
\section{PESTLE Analysis}

\textbf{From class: PESTLE = macro-environment.} Political, Economic, Social, Technological, Legal, Environmental. Here's how it applies to Tierra Bar:

\begin{description}[leftmargin=2.5em,style=nextline]
\item[Political] Sociedad Limitada, Madrid health and operating licenses, alcohol license, labour laws.
\item[Economic] Student budgets limited; tourism high (11M); inflation; unemployment 9\% (lowest since 2008).
\item[Social] Demand for casual Mexican, late-night venues, veggie/gluten-free options, craft beer. Community matters.
\item[Technological] Reviews (Google, Yelp), social media, delivery apps (Glovo, Just Eat), contactless payment.
\item[Legal] AESAN, HACCP, alcohol licensing, minimum wage (1,134 euros full-time, 620 euros part-time).
\item[Environmental] Sustainable packaging, waste reduction.
\end{description}

\vspace{1em}

% ===== 6. SWOT =====
\section{SWOT Analysis}\label{sec:swot}

\textbf{From class:} SWOT = internal (S/W) + external (O/T). Be specific.

\begin{center}
\small
\begin{tabular}{@{}p{5.2cm}p{5.2cm}@{}}
\toprule
\rowcolor{lightorange}
\textbf{Strengths} & \textbf{Weaknesses} \\
\midrule
Location (Plaza Isabel II), 9/10 reviews, customisable food, 16 craft taps, late hours, veggie options, delivery & No student programme yet, no study tables, events not established \\
\addlinespace
\rowcolor{lightorange}
\textbf{Opportunities} & \textbf{Threats} \\
\midrule
Complutense, UAM, Erasmus, tourism, expansion (Chamber\'i, Moncloa) & Taco Bell, taquer\'ias, price fluctuations, seasonal tourism, delivery commissions \\
\bottomrule
\end{tabular}
\end{center}

\vspace{1em}

% ===== 7. ECONOMIC \& SOCIAL =====
\section{Economic and Social Analysis}

Madrid received over 11 million tourists in 2025. Casual Mexican food and craft beer are in demand. Inflation raises food and alcohol costs; unemployment at 9\% (lowest since 2008) supports spending. Students and young professionals seek affordable venues. The market exists; cash follows when customers return.

\vspace{1em}

% ===== 8. COMPETITION =====
\section{Competitive Analysis}

\textbf{So who else is in this market?} From class: who serves the market? What's the gap?

Competitors: Taco Bell (chain, different positioning), local taquer\'ias (smaller, less bar-focused), general bars in Centro. Tierra Bar differentiates through customisable food, 16 craft taps, margaritas, late hours, and reviews. Almost no competitor targets students with discounts and events. Going where competitors are confirms the market; our student programme creates differentiation. Resources (location, kitchen, bar) and capabilities (events, partnerships) support this.

\vspace{1em}

% ===== 9. STRATEGY =====
\section{Strategic Plan}

\textbf{Objectives:} (1) Make Tierra Bar the Mexican bar of choice for students in Centro. (2) Retain tourists and locals. Differentiation via student focus, events, discount, Student Pass.

\textbf{Short-term (0--12 months):} Launch 15\% discount and Student Pass. Reserve study tables. First Taco Tuesday and Margarita Meetup. Flyers at Complutense and UAM. Scale events if attendance justifies.

\textbf{Long-term (1--3 years) from class:} Establish as the Mexican bar for study abroad. Second location (Chamber\'i, Moncloa). Focus: university partnerships, retention, word-of-mouth. Quote from Week 5: ``Strategy is for 15 years, not for 15 days. Slow is fast.''

\vspace{1em}

% ===== 10. CONTINGENCY =====
\section{Contingency Plan}

\textbf{From class:} Contingency plan covers key scenarios. Here's what we do if things go wrong:

\begin{itemize}[leftmargin=*,nosep]
\item \textbf{Low student uptake:} Intensify marketing at Complutense and UAM. Survey and adjust. After 12 months, refocus on tourists and locals if needed.
\item \textbf{Competitor copy:} Emphasise events, partnerships, craft beer, study tables. We have the infrastructure.
\item \textbf{Rising costs:} Negotiate suppliers; small price adjustment if necessary; keep student offer reasonable.
\item \textbf{Seasonal dip:} Focus on locals and semester Student Passes. Erasmus peaks: Sept, Feb.
\item \textbf{Events underperform:} Change format, timing, or type (e.g. quiz night, live music).
\end{itemize}

\vspace{1em}

% ===== 11. MARKETING =====
\section{Marketing Plan}

\textbf{From class:} Marketing plan needs market research, 4P's, calendar, budget Year 1. Let's break it down.

\textbf{Research:} Qualitative (student interviews: preferences, events, willingness to pay) and quantitative (survey: usage, frequency, demographics).

\textbf{4P's (from class):}
\begin{itemize}[leftmargin=*,nosep]
\item \textbf{Product:} Mexican bar, burritos/tacos/nachos, craft beer, margaritas, student discount, Student Pass, events
\item \textbf{Price:} \`a la carte; 15\% student discount; Student Pass 20 euros/month
\item \textbf{Place:} Madrid Centro, Glovo, Just Eat, walk-in
\item \textbf{Promotion:} Instagram, TikTok, Yelp, Google, university flyers, word-of-mouth
\end{itemize}

\textbf{Timeline:} M1: social media, discount, Student Pass. M2: flyers, student verification. M3: first events. M4--6: add event type, partnership talks. M7--9: Erasmus cohort, push Pass. M10--12: evaluate and plan Year 2.

\textbf{Budget Year 1:}
\begin{center}
\small
\begin{tabular}{@{}p{8cm}r@{}}
\toprule
\rowcolor{lightorange}
\textbf{Item} & \textbf{Euros} \\
\midrule
Social media & 500 \\
Flyers & 150 \\
Event materials & 250 \\
Student Pass promotion & 100 \\
University partnership & 50 \\
\midrule
\textbf{Total} & \textbf{1,050} \\
\bottomrule
\end{tabular}
\end{center}

\vspace{1em}

% ===== 12. LINK TO COURSE CONCEPTS =====
\section{Link to Course Concepts}\label{sec:venture}

\textbf{So what did we learn in class that applies here?}

\textbf{Location is the masterpiece:} Plaza Isabel II, Madrid Centro. Near Opera, tourism, students. Right place.

\textbf{Go where competitors are:} Taco Bell, taquer\'ias validate the market. Tierra Bar differentiates with student focus, events, craft beer, late hours.

\textbf{Resources vs capabilities:} Resources: kitchen, bar, 16 taps, space, ingredients. Capabilities: events, university partnerships. From the course.

\textbf{Four phases:} Conceive (plan) $\rightarrow$ Execute (improvement) $\rightarrow$ Establish (run) $\rightarrow$ Expand (scale). We're in Execute.

\textbf{Internationalization:} Madrid, Barcelona, Salamanca already. Geographic expansion follows the course's internationalization logic.

\textbf{If you take away one thing:} Business plan = opportunity + market + how to pursue. Single story. That's what we're doing here.

\vspace{1.5em}
\noindent
\begin{center}
\fcolorbox{tierraorange}{lightorange}{\parbox{0.92\textwidth}{\small\textcolor{charcoal}{\textbf{Sources:} Class notes weeks 1--7, Global Management Project.pdf, Business plan reunidas. No external sources.}}}
\end{center}

\end{document}
